% Image credits for the photographs and AI illustrations of Book 2
% (University Chemistry, Year 1), French edition. Machine-readable license
% records live in images/book2/CREDITS.md; keep this file in sync with the
% English frontmatter/image-credits-book2.tex.
\chapter*{Crédits des images}

Les dessins de ce livre sont originaux. Les photographies sont utilisées
sous leurs licences libres respectives, avec nos remerciements à leurs
auteurs~; les liens vers les sources et les adresses des licences de
chaque photographie sont recensés dans
\texttt{images/book2/CREDITS.md}, dans le dépôt du livre.
Les pictogrammes de danger sont ceux du Système général harmonisé de
classification et d'étiquetage des produits chimiques, tels que les
dessine l'extension \texttt{ghsystem} de \TeX\ Live.

\bigskip
Outre les photographies et les dessins originaux, certaines figures sont
des illustrations générées par intelligence artificielle, produites avec
le générateur d'images d'OpenAI à partir de consignes rédigées par les
éditeurs du livre, et relues pour leur exactitude chimique avant
inclusion. La consigne complète de chaque image générée est consignée
dans \texttt{images/book2/ai/PROMPTS.md}, dans le dépôt.

\bigskip
\noindent Photographies, par chapitre~:
\begin{itemize}\small
  \item Chap.~1~: Niels Bohr, vers 1922, AB Lagrelius \& Westphal,
        domaine public (Wikimedia Commons).
  \item Chap.~2~: Linus Pauling, 1962, Fondation Nobel, domaine public
        (Wikimedia Commons).
  \item Chap.~5~: échantillon de cuivre natif, par Flipin, OG, CC0
        (Wikimedia Commons).
  \item Chap.~6~: cristaux de halite, fluorine et diamant octaédrique,
        par James St. John, CC BY 2.0~; graphite, par Zbynek Burival,
        CC BY-SA 4.0~; cristaux de neige par Wilson A. Bentley (1902),
        domaine public (tous sur Wikimedia Commons).
  \item Chap.~8~: Svante Arrhenius, 1909, Meisenbach Riffarth \& Co.,
        domaine public (Wikimedia Commons).
  \item Chap.~11~: précipités de chlorure, de bromure et d'iodure
        d'argent, par Rrausch1974, CC BY-SA 3.0 (Wikimedia Commons).
  \item Chap.~12~: suspension d'hydroxyde de cuivre(II), par Alvy16,
        CC BY 4.0~; solution d'ammine de cuivre, par Chemicalinterest,
        domaine public (tous deux sur Wikimedia Commons).
  \item Chap.~13~: Walther Nernst, 1889, Smithsonian Libraries, domaine
        public (Wikimedia Commons).
  \item Chap.~16~: Jacobus Henricus van 't Hoff, photographie publiée en
        1911, domaine public (Wikimedia Commons).
  \item Chap.~21~: Victor Grignard, domaine public (Wikimedia Commons).
  \item Chap.~25~: Vladimir Markovnikov, portrait tiré d'une notice
        nécrologique publiée en 1905, domaine public (Wikimedia Commons).
  \item Chap.~26~: bassins de lithium du salar d'Atacama, NASA Earth
        Observatory, domaine public~; sodium dans l'huile minérale, par
        Justin Urgitis, CC BY-SA 2.5~; couleurs de flamme, par Hegelrast,
        CC BY-SA 4.0 (tous sur Wikimedia Commons).
  \item Chap.~27~: bauxite pisolithique, par James St. John, CC BY 2.0~;
        borax, par Aram Dulyan, domaine public~; quartz, par Stephanie
        Clifford, CC BY 2.0~; phosphore blanc, Hi-Res Images of Chemical
        Elements, CC BY 3.0~; Fritz Haber, Fondation Nobel, domaine public
        (tous sur Wikimedia Commons).
  \item Chap.~28~: extraction du soufre dans le cratère du Kawah Ijen,
        par Sémhur, CC BY-SA 4.0~; dichlore en ampoule, par W. Oelen,
        CC BY-SA 3.0~; dibrome en ampoule, par Jurii, CC BY 3.0~;
        cristaux de diiode, par Dnn87, CC BY 3.0 (tous sur Wikimedia
        Commons).
  \item Chap.~29~: banc chauffant (banc Kofler), par HBR, CC BY-SA 3.0~;
        réfractomètre d'Abbe, par Foreade, CC BY-SA 4.0 (tous deux sur
        Wikimedia Commons).
\end{itemize}
\clearpage
